\section{Comportamiento de la L-Cisteína}

La L-cisteína es un aminoácido que destaca por poseer un grupo funcional
(\ce{-SH}) en su cadena lateral, además de los grupos característicos y
generales de todos los aminoácidos, los grupos amino (\ce{-NH_2}) y carboxilo
(\ce{-COOH}). Esta estructura específica tiene ciertas particularidades que la
hacen perfecta para el diseño de biosensores, el  grupo señalado (\ce{-SH}) o
grupo tiol tiene afinidad química con los metales, en cambio, los grupos
amino y carboxilo quedan libres a interactuar con el medio o con otras
moléculas.

La concentración de la L-cisteína es un parámetro importante que dicta el
comportamiento conjunto del sistema. En bajas concentraciones contribuye a
distribuir uniformemente sobre la superficie de las nanopartículas,
manteniendo  una repulsión electrostática que preserva la estabilidad de la
misma. A este sistema se le denomina ``modo estable'' y es un requisito
fundamental para la configuración del biosensor. Sin embargo, si la
concentración es relativamente alta, los grupos funcionales libres interactúan
entre sí induciendo una agregación  de las nanopartículas sin la necesidad de
un analito externo, debido a dichas interacciones se le da el nombre de
``modo inestable''.

\subsection{Enlace entre las AuNP's y la L-Cisteína}

Durante la síntesis, las nanopartículas de oro se estabilizan con iones de
citrato derivados del citrato trisodio, esto previene su aglomeración entre
ellos. En la fase de funcionalización estas nanopartículas se exponen a la
L-cisteína. El proceso se fundamenta en la formación de un enlace covalente
entre los átomos de oro (\ce{Au}) y el átomo de azufre (\ce{S}) del grupo tiol
de la cisteína.

El enlace \ce{Au-S} es termodinámicamente mucho más estable que las interacciones
con el citrato provocando que la cisteína desplace a las moléculas
estabilizadoras y modifique la capa dieléctrica de la nanopartícula.
Al realizar este reemplazo con una concentración controlada de
\qty{0.1}{\milli\molar} de cisteína, el pico de resonancia de plasmón
superficial se mantiene prácticamente constante, indicando que las propiedades
ópticas de las nanopartículas individuales se conservan para la posterior etapa
de sensado.

\subsection{Detección de Arsénico}
La capacidad de las AuNPs funcionalizadas con cisteína (AuNPs+Cis) para detectar
iones de arsénico (\ce{As^{3+}}) se basa en la química de coordinación y en las
propiedades ópticas (efecto plasmónico) de las nanopartículas. Cuando se
introduce el \ce{As^{3+}} en la dispersión, este ion actúa como una molécula que
busca estabilizarse atrayendo densidad electrónica.

Los grupos carboxilo (\ce{-COOH}) y amino libres de la L-cisteína anclada actúan
como agentes quelantes. Un solo ion de \ce{As^{3+}} tiene la capacidad de
coordinarse simultáneamente con los grupos funcionales de cisteínas que están
unidas a diferentes nanopartículas de oro, actuando como un puente molecular
que las une, esta formación  provoca una agregación rápida de las AuNPs.

A medida que las nanopartículas se acercan debido a este puente iónico, ocurre
un acoplamiento de sus campos electromagnéticos superficiales, alterando
fuertemente las propiedades dieléctricas de su entorno local. Ópticamente, esto
se manifiesta como un efecto batocrómico (un corrimiento de la banda de absorción
hacia longitudes de onda más largas). Visualmente, la dispersión coloidal cambia
de su característico color rojo a tonos púrpuras, un cambio perceptible a simple
vista (efecto hipocrómico y batocrómico). Esta agregación mediada por el arsénico
es altamente sensible y depende del tamaño de las AuNPs, permitiendo un límite
mínimo de detección de \qty{0.07}{ppm}, un valor muy cercano al límite máximo de
\qty{0.01}{ppm} recomendado por la Organización Mundial de la Salud para agua
potable.
